% Post-checkpoint 07 extension.  This fragment is separate pending independent audit.
\begin{proposition}[Proved: a Q-factorial extraction for a formally graded surface germ]
\label{D:weighted-graded-extraction}
Let $k$ be algebraically closed of characteristic $p>0$, and let
$(A,\mathfrak m)$ be an excellent normal local $k$-algebra of
dimension two with residue field $k$.  Suppose that its completion
is isomorphic, as a complete local $k$-algebra, to the completion
at $R_+$ of a finitely generated normal positively graded domain
\[
 R=\bigoplus_{d\geq0}R_d,\qquad R_0=k,\qquad \dim R=2.
\]
There is a projective birational morphism
$h:Y\to\operatorname{Spec}A$, an isomorphism on the punctured
spectrum, with the following properties.  The surface $Y$ is
normal and locally Q-factorial, its exceptional locus has exactly
one irreducible curve, and every finite cover of the punctured
spectrum is tame at the valuation of that curve.  The exceptional
valuation, after completion at $\mathfrak m$, is the radial
valuation of the grading.
\end{proposition}
\begin{proof}
We divide all occurring degrees by their greatest common divisor,
so the grading is effective.  Choose $N>0$ such that the Veronese
ring $R^{(N)}$ is generated in degree one, using
\cite[Lemma 10.56.2, tag 0EGH]{C_StacksWeightedCone}.
Put $S=R^{(N)}$ and $C=\operatorname{Proj}S=\operatorname{Proj}R$.
The ring $S$ is normal, $R$ is finite over $S$, and $C$ is a smooth
projective integral curve.  Here normality of the Veronese follows
by intersecting $R$ with its degree-multiple fraction subfield;
normality of $C$ follows from its degree-zero affine charts.

We first construct the model over $R$.  The blow-up
\[
 W=\operatorname{Bl}_{S_+}(\operatorname{Spec}S)
\]
is the total space of $\mathcal O_C(-1)$.  Indeed, for a nonzero
$f\in S_1$ its blow-up chart has ring
\[
 S[S_1/f]=(S[f^{-1}])_0[f].
\]
The element $f$ is algebraically independent over the degree-zero
fraction field, and the changes between these charts are the
linear changes of a line-bundle coordinate.  Write
$F=k(C)$ and $K=\operatorname{Frac}R$.  Effectivity gives a rational
homogeneous element $t$ of degree one, by taking a product of
homogeneous elements with integer exponents.  Thus
\[
 K=F(t),\qquad \operatorname{Frac}S=F(t^N).
\]
Let $Y_R$ be the normalization of $W$ in $K$.  It is finite over
$W$.  The elements of $R$ are integral over $S$, hence regular
on $Y_R$, so this construction gives a projective birational
morphism $Y_R\to\operatorname{Spec}R$, an isomorphism away from
the vertex.  Projectivity follows also by factoring the projective
morphism $Y_R\to\operatorname{Spec}S$ through the separated finite
morphism $\operatorname{Spec}R\to\operatorname{Spec}S$.

At the generic point of the zero section of $W$, a line coordinate
$u$ satisfies $u=f t^N$ for some $f\in F^*$.  The extension
of the $u$-adic valuation to $F(t)$ is totally ramified of degree
$N$, with residue field $F$, and is unique.  This follows directly
from the Eisenstein polynomial $T^N-u/f$ over $F[[u]]$, including
when $p$ divides $N$.  Consequently $Y_R$ has exactly one
exceptional curve.  Its valuation is the degree valuation of $R$.

We verify local Q-factoriality, retaining the characteristic-$p$
possibilities in this construction.  At $c\in C$, choose a
uniformizer $z$ and write $f=z^b\epsilon$ with $\epsilon$ a unit.
Absorb $\epsilon$ into $u$ and multiply $t$ by a suitable power of
$z$.  The local finite normalization is then given by normalizing
\begin{equation}\label{D:weighted-toric-chart}
 T^N=u z^a,\qquad 0\leq a<N.
\end{equation}
We may take an affine neighborhood of $c$ on which $z$ defines an
\'etale morphism to $\mathbb A^1_k$.  Normalization commutes with
this \'etale base change
\cite[Lemma 37.19.2, tag 03GV]{C_StacksWeightedCone}.
The normalization of the displayed affine equation over
$k[u,z]$ has a finite surjection from $\mathbb A^2_k$, given by
\[
 u=r^N,\qquad z=s^N,\qquad T=r s^a.
\]
Indeed, $r$ and $s$ are integral, the polynomial ring $k[r,s]$
is normal, and its field extension over $k(u,z,T)$ has degree
$N$.  To compute the degree, first
$[k(r,s):k(r^N,s^N)]=N^2$, whereas adjoining $r s^a$ to
$k(r^N,s^N)$ has degree $N$, as its $r$-valuation has value one.
The resulting finite map factors through the normalization.
After the \'etale base change in $z$, its source remains regular.

For completeness, a normal local domain admitting a finite
surjection from a regular semilocal scheme is Q-factorial.  Pull
a Weil divisor $D$ back using the ramification indices of the
height-one valuations.  On the regular semilocal source this
divisor is principal.  Taking the product of the field norms on
its components and pushing the divisor down makes $mD$ principal,
where $m$ is the sum of their generic degrees.  The identities
used here are
\[
 f_*f^*D=mD,\qquad
 f_*\operatorname{div}(g)=\operatorname{div}(\operatorname{Nm}(g));
\]
the second is the norm formula for principal divisors
\cite[Lemma 42.18.1, tag 02RS]{C_StacksWeightedCone}, and the first
is the degree formula for a finite extension of DVRs.  Regular
semilocal schemes have trivial Picard groups and are locally
factorial \cite[tags 02M9 and 0AG0]{C_StacksWeightedCone}.
This argument does not require a separable field
extension.  Applied to our local charts it shows that $N$ kills
each local divisor class along the exceptional curve.  The same
argument holds after completing at any of these points: the
finite cover completes to a product of regular complete local
rings.  Thus these completed local rings are Q-factorial as well.

The model just constructed is a normalized blow-up of a
vertex-primary ideal.  More explicitly, put
\[
 J=R_N R=S_+R.
\]
The equality uses that $S$ is generated in degree one, and $J$
is $R_+$-primary because $R$ is integral and graded over $S$.
The blow-up of $J$ is the dominant component of the base change
of $W$ to $R$: the Rees algebra map
\[
 \left(\bigoplus_{n\geq0}S_+^n\right)\otimes_S R
 \longrightarrow \bigoplus_{n\geq0}J^n
\]
is surjective and is an isomorphism on the dense complement of
the vertex.  In particular, this blow-up is finite over $W$.
Its normalization is therefore $Y_R$, by uniqueness of the
normalization of $W$ in $K$.

Identify $\widehat A$ with the completed graded local ring.
The ideal $J\widehat A$ contains a power of the maximal ideal,
so it is the extension of an ideal $I\subset A$: use
$A/\mathfrak m^q=\widehat A/\mathfrak m^q\widehat A$ for a
sufficiently large $q$.  Let $Y$ be the normalization of
$\operatorname{Bl}_I(\operatorname{Spec}A)$.  Its base change to
$\widehat A$ is the normalization of the corresponding completed
blow-up.  This follows directly from
\cite[Lemmas 54.11.3 and 54.11.6, tag 0BG4]{C_StacksWeightedCone}:
away from the special fibre the blow-up is the normal punctured
spectrum.  In particular, the normalization is finite and $Y$
is projective over $A$.

There is exactly one exceptional curve, since the special fibres
before and after completion are identical.  To descend
Q-factoriality, one can work with the completed local rings at
points of this curve.  These rings are unchanged by base change
to $\widehat A$
\cite[Lemma 54.11.1, tag 0BG4]{C_StacksWeightedCone} and, by the
finite regular charts above, are Q-factorial.  If a reflexive
rank-one module $M$ over a normal local ring has an invertible
reflexive power after completion, that power is already
invertible.  Indeed, finite-module duals and double duals commute
with flat base change
\cite[Lemma 17.22.5, tag 01CM]{C_StacksWeightedCone},
and finite locally free modules descend
faithfully flatly.  This proves local Q-factoriality on $Y$;
outside the exceptional fibre its local rings are regular.
The normality needed after the intermediate completion base
changes follows from regularity of completion maps of excellent
rings and normality under regular base change
\cite[tags 0BFK and 0H7S]{C_StacksWeightedCone}.

Finally, pull any finite punctured cover back to $\widehat A$ and
take the connected components and their Galois closures.
Proposition~\ref{D:graded-cone-tame} makes them tame at the
radial valuation.  The original exceptional DVR and the radial
DVR on the completed model have the same actual completion, by
Lemma 54.11.1 applied at the generic point of the special-fibre
curve.  Completion preserves the ramification index and the
separability of the residual extension for finite extensions of
excellent DVRs.  Tameness therefore descends to the original
exceptional valuation, as required.
\end{proof}

\begin{corollary}[Proved: formally graded points in a proper surface boundary]
\label{D:weighted-graded-boundary}
In Theorem~\ref{D:formal-cone-mixed-boundary} and
Corollary~\ref{D:formal-cone-fixed-boundary}, we may add the
following alternative for every omitted non-Q-factorial point:
its completed local ring is a positively graded normal surface
completion as in Proposition~\ref{D:weighted-graded-extraction}.
Thus the punctured fundamental group is topologically finitely
generated in the finite-boundary case, and the fixed-boundary
tame quotient is topologically finitely generated in the general
boundary case with these non-Q-factorial points isolated in the
boundary.  The finite-normalization version of the first theorem
has the same additional alternative on the normalizations.
\end{corollary}
\begin{proof}
At each formally graded omitted point use the preceding
projective partial resolution.  It spreads to an affine
neighborhood by finite presentation.  The ideal may be extended
to a coherent ideal equal to the unit ideal off that point, so
the normalized blow-up is an isomorphism away from the prescribed
point.  Keep the omitted Q-factorial points and resolve the
retained singularities as in the cited mixed-boundary arguments.
The resulting proper normal surface is locally Q-factorial and
hence projective by Lemma~\ref{D:Qfactorial-partial-resolution}.
At a formally graded point the only new boundary divisor is its
radial curve, and all the covers under consideration are tame
there by Proposition~\ref{D:weighted-graded-extraction}.  No
divisors arise from the Q-factorial singularities retained on
the partial resolution.  The remaining boundary divisors are
exactly those already handled in the mixed-boundary theorem.
The projective fixed-boundary tame group consequently surjects
onto the required group.  Its finite generation gives the claim.
The finite-normalization descent is unchanged.
\end{proof}
